%%%%%%%%%%%%%%%%%%%%%%%%%%%%%%%%%%%%%%%%%%%%%%%%%%%%%%%%%%%%%%%%%%%%%%%%%%%%%%%%
%% APPENDIX A: SUPPORTING MATERIAL FOR CHAPTER 1 (INTRODUCTION)
%% Research Traceability Matrix
%%%%%%%%%%%%%%%%%%%%%%%%%%%%%%%%%%%%%%%%%%%%%%%%%%%%%%%%%%%%%%%%%%%%%%%%%%%%%%%%

\chapter{Research Traceability Matrix (Chapter~1)}
\label{app:ch1_support}
\label{app:traceability}

This appendix provides the complete traceability matrix linking the business problem (Chapter~1) through literature gaps (Chapter~2), research design (Chapter~3), analysis and findings (Chapter~4), and conclusions (Chapter~5). This matrix demonstrates the logical inevitability and rigor of the research, ensuring that every design decision traces back to an identified gap and every finding traces forward to a validated contribution.

\section{Literature Theme to Gap Traceability}

\noindent Table~\ref{tab:trace_t1} maps each literature theme to the specific gap it leaves unaddressed and the corresponding DBA relevance.
\needspace{8\baselineskip}
{\tablestripes\tablefontsize
\begin{xltabular}{\textwidth}{@{} L L L L @{}}
\caption{Literature-to-Gap Traceability (T1)}\label{tab:trace_t1} \\
\toprule
\thead{Literature Theme} & \thead{What Literature Covers} & \thead{What Is Missing (Gap)} & \thead{DBA Relevance} \\
\midrule
\endfirsthead
\multicolumn{4}{@{}l}{\textit{(\,Tab.~\ref{tab:trace_t1} continued from previous page\,)}} \\
\toprule
\thead{Literature Theme} & \thead{What Literature Covers} & \thead{What Is Missing (Gap)} & \thead{DBA Relevance} \\
\midrule
\endhead
\bottomrule
\endlastfoot
AI diagnostic systems & Algorithm accuracy in single domains & ASD-focused unified pipeline & Scalability and cost efficiency \\
\addlinespace
AI governance & Ethical principles and policies & Operational governance execution & Risk and accountability \\
\addlinespace
Clinical decision support & Technology-centric tools & Decision ownership and human-AI interaction & Managerial clarity \\
\addlinespace
Explainable AI & SHAP, Grad-CAM methods & Literature-grounded clinical explanations & Trust and regulatory readiness \\
\addlinespace
Workflow automation & Task-specific automation & End-to-end agentic orchestration & Efficiency and reproducibility \\
\end{xltabular}}
\vspace{0.5em}\noindent\textit{Note.} AI = artificial intelligence; SHAP = SHapley Additive exPlanations; DBA = Doctor of Business Administration.

\section{Gap to Research Objective Mapping}

\noindent Table~\ref{tab:trace_t2} traces each identified gap to its corresponding research objective and the intended DBA outcome.
\needspace{8\baselineskip}
{\tablestripes\tablefontsize
\begin{xltabular}{\textwidth}{@{} L L L @{}}
\caption{Gap to Research Objectives Traceability (T2)}\label{tab:trace_t2} \\
\toprule
\thead{Identified Gap} & \thead{Research Objective} & \thead{DBA Intent} \\
\midrule
\endfirsthead
\multicolumn{3}{@{}l}{\textit{(\,Tab.~\ref{tab:trace_t2} continued from previous page\,)}} \\
\toprule
\thead{Identified Gap} & \thead{Research Objective} & \thead{DBA Intent} \\
\midrule
\endhead
\bottomrule
\endlastfoot
No unified ASD-focused pipeline & Design unified AI framework across the ASD domain & Scalable diagnostic capability \\
\addlinespace
Lack of decision ownership & Define decision-centric governance model & Managerial accountability \\
\addlinespace
Governance not operationalized & Embed controls into AI workflows & Risk reduction \\
\addlinespace
AI value not measurable & Define KPI and ROI measurement model & Executive relevance \\
\addlinespace
Human--AI interaction ambiguous & Formalize interaction model with clear roles & Trust and adoption \\
\addlinespace
Limited explainability & Integrate RAG for literature-grounded explanations & Regulatory compliance \\
\end{xltabular}}
\vspace{0.5em}\noindent\textit{Note.} AI = artificial intelligence; KPI = key performance indicator; ROI = return on investment; RAG = retrieval-augmented generation.

\section{Research Objective to Framework Component Mapping}

\noindent Table~\ref{tab:trace_t3} links each research objective to its framework element, Chapter~4 artifact, and supporting evidence.
\needspace{8\baselineskip}
{\tablestripes\tablefontsize
\begin{xltabular}{\textwidth}{@{} L L L L @{}}
\caption{Objectives to Framework Components (T3)}\label{tab:trace_t3} \\
\toprule
\thead{Research Objective} & \thead{Framework Element} & \thead{Chapter~4 Artifact} & \thead{Evidence} \\
\midrule
\endfirsthead
\multicolumn{4}{@{}l}{\textit{(\,Tab.~\ref{tab:trace_t3} continued from previous page\,)}} \\
\toprule
\thead{Research Objective} & \thead{Framework Element} & \thead{Chapter~4 Artifact} & \thead{Evidence} \\
\midrule
\endhead
\bottomrule
\endlastfoot
Unified diagnostic pipeline & Five-layer architecture & Framework Section & Architecture diagram \\
\addlinespace
Decision ownership & Decision classification model & Decision Matrix & Decision matrix \\
\addlinespace
Operationalize governance & Governance controls matrix & Governance Section & Control inventory \\
\addlinespace
Human--AI interaction & Decision sequence model & Interaction Model & Sequence diagram \\
\addlinespace
Measure value & KPI and ROI model & KPI Section & Metrics framework \\
\addlinespace
Roles and accountability & RACI matrix & RACI Section & Role definitions \\
\end{xltabular}}
\vspace{0.5em}\noindent\textit{Note.} RACI = responsible, accountable, consulted, informed; KPI = key performance indicator; ROI = return on investment.

\section{Comparative Justification Matrix}

\noindent Table~\ref{tab:trace_t4} compares five existing governance and AI models against the proposed framework, identifying each model's limitation and the corresponding remedy.
\needspace{8\baselineskip}
{\tablestripes\tablefontsize
\begin{xltabular}{\textwidth}{@{} L L L L @{}}
\caption{Existing Models vs.\ This Framework (T4)}\label{tab:trace_t4} \\
\toprule
\thead{Existing Model} & \thead{Strength} & \thead{Limitation} & \thead{How Framework Fixes It} \\
\midrule
\endfirsthead
\multicolumn{4}{@{}l}{\textit{(\,Tab.~\ref{tab:trace_t4} continued from previous page\,)}} \\
\toprule
\thead{Existing Model} & \thead{Strength} & \thead{Limitation} & \thead{How Framework Fixes It} \\
\midrule
\endhead
\bottomrule
\endlastfoot
AI capability models & Technical depth & No accountability model & Decision ownership layer \\
\addlinespace
Ethics guidelines & Normative clarity & Non-operational & Embedded governance controls \\
\addlinespace
IT governance (COBIT) & Process rigor & AI-agnostic & AI-specific risk controls \\
\addlinespace
Maturity models & Assessment structure & No execution pathway & Operational workflow \\
\addlinespace
Single-domain AI & High accuracy & Not generalizable & ASD-focused unified pipeline \\
\end{xltabular}}
\vspace{0.5em}\noindent\textit{Note.} COBIT = Control Objectives for Information and Related Technologies; AI = artificial intelligence.

\section{Research Question to Framework Evidence}

\noindent Table~\ref{tab:trace_t5} maps each research question to the chapter section where it is answered and the type of evidence provided.
\needspace{8\baselineskip}
{\tablestripes\tablefontsize
\begin{xltabular}{\textwidth}{@{} C L L L @{}}
\caption{Research Questions to Framework Evidence (T5)}\label{tab:trace_t5} \\
\toprule
\thead{RQ} & \thead{Question} & \thead{Where Answered} & \thead{Evidence Type} \\
\midrule
\endfirsthead
\multicolumn{4}{@{}l}{\textit{(\,Tab.~\ref{tab:trace_t5} continued from previous page\,)}} \\
\toprule
\thead{RQ} & \thead{Question} & \thead{Where Answered} & \thead{Evidence Type} \\
\midrule
\endhead
\bottomrule
\endlastfoot
1 & Can unified pipeline work across domains? & Ch.~4 & Performance tables \\
\addlinespace
2 & Does DL outperform classical baselines? & Ch.~4 & Statistical tests \\
\addlinespace
3 & What features contribute most? & Ch.~4 & SHAP analysis \\
\addlinespace
4 & Can agentic AI automate workflows? & Ch.~4 & Ablation study \\
\addlinespace
5 & Does RAG improve explainability? & Ch.~4 & Expert ratings \\
\addlinespace
6 & What ASD-focused patterns exist? & Ch.~4 & Feature heatmap \\
\end{xltabular}}
\vspace{0.5em}\noindent\textit{Note.} RQ = research question; DL = deep learning; SHAP = SHapley Additive exPlanations; RAG = retrieval-augmented generation.

\section{Framework to Validation to Value}

\noindent Table~\ref{tab:trace_t6} documents the complete chain from framework element through validation method to business outcome.
\needspace{8\baselineskip}
{\tablestripes\tablefontsize
\begin{xltabular}{\textwidth}{@{} L L L L @{}}
\caption{Framework Outcome Traceability (T6)}\label{tab:trace_t6} \\
\toprule
\thead{Framework Element} & \thead{Validation Method} & \thead{Business Outcome} & \thead{Chapter} \\
\midrule
\endfirsthead
\multicolumn{4}{@{}l}{\textit{(\,Tab.~\ref{tab:trace_t6} continued from previous page\,)}} \\
\toprule
\thead{Framework Element} & \thead{Validation Method} & \thead{Business Outcome} & \thead{Chapter} \\
\midrule
\endhead
\bottomrule
\endlastfoot
Decision classification & Expert review & Reduced decision ambiguity & Ch.~4 \\
\addlinespace
Governance controls & Scenario testing & Risk mitigation & Ch.~4 \\
\addlinespace
Operational workflow & Practitioner feedback & Adoption readiness & Ch.~4 \\
\addlinespace
KPI model & Before/after comparison & ROI visibility & Ch.~4, Ch.~5 \\
\addlinespace
Human--AI sequence & Case application & Trust and accountability & Ch.~4 \\
\addlinespace
Five-layer architecture & ASD-focused validation & Scalable deployment & Ch.~4, Ch.~5 \\
\end{xltabular}}
\vspace{0.5em}\noindent\textit{Note.} KPI = key performance indicator; ROI = return on investment; AI = artificial intelligence.
